%=============================================================================
% MODULE 03: INTRODUCTION
%=============================================================================
% Frontiers in Public Health — Funnel structure:
%   broad context → problem → gap → objectives → contribution
%=============================================================================

\section{Introduction}

\subsection{Background}

Artificial intelligence (AI) and machine learning (ML) are transforming 
healthcare delivery across multiple domains, from diagnostic imaging to 
clinical decision support systems \cite{topol2019}. In diabetes care 
specifically, ML models demonstrate substantial promise for early detection, 
risk stratification, and personalized treatment recommendations 
\cite{sd2019, esteva2017}. The global diabetes epidemic, affecting over 
537 million adults worldwide \cite{idf2021}, demands scalable and accurate 
screening tools that can operate effectively across diverse populations.

However, growing evidence suggests that these technologies may perpetuate 
or amplify existing health disparities through algorithmic bias 
\cite{mehrabi2021}. Algorithmic bias in healthcare refers to systematic 
and unfair discrimination against certain groups, often rooted in historical 
data that reflects societal inequities \cite{barocas2019}. In diabetes 
diagnosis, biased models may produce less accurate predictions for 
underrepresented populations, leading to delayed diagnoses, inappropriate 
treatments, and ultimately worse health outcomes \cite{obermeyer2019}.

The diabetes epidemic disproportionately affects marginalized communities. 
In the United States, diabetes prevalence is 60\% higher among Black adults 
and 77\% higher among Hispanic adults compared to White adults \cite{cdc2020}. 
These disparities are driven by social determinants of health, including 
access to healthcare, nutrition, and socioeconomic factors \cite{hill2022}. 
When ML models are trained on data that underrepresents these populations, 
they may fail to generalize effectively, exacerbating existing inequities.

\subsection{Problem Statement}

Despite growing awareness of algorithmic fairness, few studies have 
systematically evaluated bias in diabetes prediction models across multiple 
demographic dimensions simultaneously. Most existing research focuses on 
single attributes (race \textit{or} gender) rather than examining 
intersectional effects where multiple forms of disadvantage converge 
\cite{crenshaw1989, buolamwini2018}. Furthermore, while technical metrics 
for fairness have been proposed \cite{corbett2018}, their practical 
implications for clinical decision-making remain insufficiently understood.

The intersectionality framework, originally articulated by Crenshaw 
\cite{crenshaw1989}, highlights how overlapping identities create unique 
modes of discrimination that cannot be captured by analyzing each attribute 
in isolation. In the context of algorithmic fairness, this means that bias 
experienced by individuals at the intersection of multiple marginalized 
groups (e.g., minority women) may exceed the sum of biases attributed to 
each individual group.

\subsection{Objectives}

This study aims to:

\begin{enumerate}[leftmargin=*]
    \item Evaluate the discriminative performance of four ML models for 
          diabetes diagnosis across demographic subgroups defined by gender 
          and ethnicity.
    \item Quantify algorithmic bias using established fairness metrics, 
          including Equalized Odds and Demographic Parity.
    \item Perform intersectional analysis to identify amplified disparities 
          at the intersection of race and gender.
    \item Provide evidence-based recommendations for bias mitigation in 
          clinical ML systems deployed in diabetes care.
\end{enumerate}

\subsection{Significance and Contribution}

This research contributes to the growing literature on healthcare AI equity 
through:

\begin{itemize}[leftmargin=*]
    \item A comprehensive evaluation framework for assessing algorithmic bias 
          that integrates performance, fairness, and intersectional analysis.
    \item Empirical evidence of intersectional disparities in diabetes 
          prediction, extending beyond single-attribute bias assessment.
    \item Practical guidance for developing fairer clinical AI systems, 
          including actionable recommendations for data collection, model 
          development, and deployment monitoring.
    \item A fully reproducible methodology with open-source code, facilitating 
          replication and extension by the research community.
\end{itemize}
